\section{Analysis of Model-\texorpdfstring{$C_2$}{C2} (\texorpdfstring{$\theta_1>0,\,\gamma_1=\gamma_2=\theta_2=0$}{theta1>0, gamma1=gamma2=theta2=0})}\label{sec:model_c2}

In this section we derive $P(x,y)$ for the model in which only class-$1$ customers abandonment, but no class-$2$ abandonment and no jockeying. Figure~\ref{fig:model-c2-queue} shows the queue layout.
\begin{figure}[H]
    \centering
    \resizebox{0.62\textwidth}{!}{%
        % Model C₂: γ₁=γ₂=θ₂=0, θ₁>0
% Class-1 customers abandon the queue at rate θ₁ (dotted, upward from Queue 1).
% No jockeying. Class-2 never abandons.
\begin{tikzpicture}[>=Latex, thick]

    % ── Queue 1 (top, Class-1, priority) ──────────────────────────────────────
    \draw (0, 1.9) -- (2, 1.9) -- (2, 1.1) -- (0, 1.1);
    \foreach \x in {0.2, 0.4, ..., 1.8} { \draw[thin] (\x, 1.8) -- (\x, 1.2); }
    \draw[->] (-1.5, 1.5) -- (-0.1, 1.5) node[midway, above] {$\lambda_1$};
    % θ₁: Class-1 abandonment (dotted, upward)
    \draw[->, dotted] (1, 2.0) -- (1, 3.0) node[above] {$\theta_1$};

    % ── Queue 2 (bottom, Class-2) ──────────────────────────────────────────────
    \draw (0, -1.1) -- (2, -1.1) -- (2, -1.9) -- (0, -1.9);
    \foreach \x in {0.2, 0.4, ..., 1.8} { \draw[thin] (\x, -1.2) -- (\x, -1.8); }
    \draw[->] (-1.5, -1.5) -- (-0.1, -1.5) node[midway, above] {$\lambda_2$};
    % No θ₂: Class-2 does not abandon

    % ── Server ────────────────────────────────────────────────────────────────
    \node[draw, circle, minimum size=1.2cm] (server) at (5, 0) {$\mu$};
    \draw[->] (2.2,  1.5) -- (server);
    \draw[->] (2.2, -1.5) -- (server);
    \draw[->] (server) -- (7, 0);

\end{tikzpicture}
%
    }
    \caption{Queue layout for Model-$C_2$. Class-1 customers abandon the system
    at rate $\theta_1$ per waiting customer (dotted arrow, upward from Queue~1).
    Class-2 never abandons ($\theta_2=0$). There is no jockeying.}
    \label{fig:model-c2-queue}
\end{figure}
This choice of parameters is motivated by two key aspects of the resulting model: first, class-$2$ still sees Poisson arrivals, so the probabilistic argument of Model-$A$ carries over; second, substituting these values into the general fundamental equation~\eqref{eq:gen:fundamental_equation} leaves a first-order linear ODE in $x$, with no $y$-derivatives the boundary function $P_y(y)$ comes out of the integral, sparing us the Volterra equation that left Model-$B$ open (\S\ref{sec:model_b}).

\textbf{Stability.} Unlike the models without abandonment, stability here is non-trivial. The class-$1$ queue is an $M/M/1{+}M$ (Erlang-A) subsystem, which is positive recurrent for \emph{every} load (\S\ref{ssec:erlangA}); The only constraint comes from class-$2$, which —explained in more detail in \S\ref{ssec:C2:prob}— sees an $M/G/1$ queue whose service time is one class-$1$ busy period $B_C$. That queue is stable exactly when its utilisation stays below one:
\begin{equation}
    \lambda_2\,\mathbb{E}[B_C] \;<\; 1
    \qquad\Longleftrightarrow\qquad
    \rho_2\cdot{}_1F_1\!\left(1;\,\tfrac{\mu}{\theta_1}+1;\,\tfrac{\lambda_1}{\theta_1}\right) \;<\; 1,
    \label{eq:C2:stability}
\end{equation}
which follows from the mean busy period $\mathbb{E}[B_C]$ (Equation~\eqref{eq:erlangA:bp_mean}) together with the Kummer identity $\mu\bigl({}_1F_1(1;\tfrac{\mu}{\theta_1};\tfrac{\lambda_1}{\theta_1})-1\bigr)=\lambda_1\,{}_1F_1(1;\tfrac{\mu}{\theta_1}+1;\tfrac{\lambda_1}{\theta_1})$. Under the stability condition (Equation~\eqref{eq:C2:stability}), the PGF exists.

The following theorem states the result, obtained by both an analytical and a probabilistic methods:
\begin{theorem}\label{thm:C2:PGF}
    Consider the two-class non-preemptive priority system of Model-$C_2$ ($\theta_1>0$, $\gamma_1=\gamma_2=\theta_2=0$): Poisson arrivals at rates $\lambda_1,\lambda_2$, common exponential service at rate $\mu$, and per-customer abandonment of \emph{waiting} class-$1$ customers at rate $\theta_1$. Under the stability condition~\eqref{eq:C2:stability}, the probability generating function is
    \begin{align*}
        \begin{split}
            P(x,y) \;&=\; \frac{e^{\lambda_1 x/\theta_1}\,\mu\,\pi(0,0)\,(1-y)}{\theta_1\, x^{\mu/\theta_1}(1-x)^{\lambda_2(1-y)/\theta_1}\bigl(\widetilde{B}_C(\lambda_2(1-y))-y\bigr)}\\[4pt]
            &\quad\times \Bigl[\widetilde{B}_C(\lambda_2(1-y))\,\mathcal{H}_1(x,y) - \bigl(1-\widetilde{B}_C(\lambda_2(1-y))\bigr)\,\mathcal{H}_2(x,y)\Bigr],
        \end{split}
    \end{align*}
    where
    \begin{align*}
        \mathcal{H}_1(x,y) &\;=\; \int_0^x e^{-\lambda_1 t/\theta_1}\,t^{\,\mu/\theta_1-1}(1-t)^{\,\lambda_2(1-y)/\theta_1}\,dt, \\
        \mathcal{H}_2(x,y) &\;=\; \int_0^x e^{-\lambda_1 t/\theta_1}\,t^{\,\mu/\theta_1}(1-t)^{\,\lambda_2(1-y)/\theta_1-1}\,dt,
    \end{align*}
    $\widetilde{B}_C(s)$ is the LST of the busy period of the class-$1$ $M/M/1{+}M$ subsystem given by~\eqref{eq:erlangA:bp_lst}, and $\pi(0,0)$ is given by Corollary~\ref{cor:C2:pi00} below.
\end{theorem}

\begin{proof}
The proof consists of three steps:

\emph{Step~(i): reduction to a first-order ODE in $x$.}  Substituting the Model-$C_2$ parameters into the general fundamental equation (Equation~\eqref{eq:gen:fundamental_equation}) leaves, for each fixed $y$, a first-order linear ODE in $x$~\eqref{eq:C2:fundamental}. 
\begin{align}
    \begin{split}
            &\left[ (\lambda_1+\lambda_2+\mu)xy - \mu y - \lambda_1 x^2 y - \lambda_2 x y^2 \right] P(x,y)
        + \theta_1\, xy(x-1)\,\frac{\partial}{\partial x} P(x,y) \\
        & \quad=\quad \mu (x-y)\,P_y(y) - \mu x(1-y)\,\pi(0,0),
        \label{eq:C2:fundamental}
    \end{split} 
\end{align}
%TODO: Add brief Partial Fraction Decomposition, aligned with notation in Model-$B$

Its integrating factor (\S\ref{sssec:if}) is
\begin{equation*}
    \mu_I(x;y) \;=\; e^{-\lambda_1 x/\theta_1}\,x^{\alpha}\,(1-x)^{\beta(y)},
    \qquad \alpha=\tfrac{\mu}{\theta_1},\quad \beta(y)=\tfrac{\lambda_2(1-y)}{\theta_1},
\end{equation*}
which vanishes at $x=0$ and, for $y<1$, at $x=1$, since $\alpha>0$ and $\beta(y)\ge0$.

\emph{Step~(ii): determination of $P_y(y)$.}  We obtain $P_y(y)$ two independent ways.  \emph{Analytically} (\S\ref{ssec:C2:anal}), boundedness of $P(x,y)$ as $x\to1$ forces $\int_0^1 q(t;y)\,\mu_I(t;y)\,dt=0$~\eqref{eq:C2:boundedness}, which the Euler--Kummer identity~\eqref{eq:C2:Bc_Kummer} solves to give $P_y(y)$~\eqref{eq:C2:Py}.  \emph{Probabilistically} (\S\ref{ssec:C2:prob}), class-$2$ sees an $M/G/1$ queue with service time $B_C$, and the Pollaczek--Khinchine formula returns the same $P_y(y)$.

\emph{Step~(iii): recovery of $P(x,y)$.}  With $P_y(y)$ known, integrating the ODE against $\mu_I$ (\S\ref{ssec:C2:recovery}) yields the closed form~\eqref{eq:C2:Pxy}.
\end{proof}

\begin{corollary}\label{cor:C2:pi00}
    The limiting probability $\pi(0,0)$ of the busy state with both queues empty, and the empty-system probability $\pi_0$, are
    \begin{align}
        \pi(0,0) \;&=\; \frac{(\lambda_1+\lambda_2)\bigl(1 - \lambda_2\,\mathbb{E}[B_C]\bigr)}{\mu\,\bigl(1 + \lambda_1\,\mathbb{E}[B_C]\bigr)} \;=\; (\rho_1+\rho_2)\,\frac{1 - \lambda_2\,\mathbb{E}[B_C]}{1 + \lambda_1\,\mathbb{E}[B_C]}, \label{eq:C2:pi00}\\[4pt]
        \pi_0 \;&=\; \frac{1 - \lambda_2\,\mathbb{E}[B_C]}{1 + \lambda_1\,\mathbb{E}[B_C]}, \label{eq:C2:pi0}
    \end{align}
    where $\mathbb{E}[B_C]$ is the mean class-$1$ busy period~\eqref{eq:erlangA:bp_mean}. As $\theta_1\to 0^+$, $\mathbb{E}[B_C]\to (\mu-\lambda_1)^{-1}$ and these recover the Model-$A$ baseline $\pi(0,0) = \rho(1-\rho)$, $\pi_0 = 1-\rho$.
\end{corollary}

\subsection{Reduced fundamental equation}\label{ssec:C2:reduction}

Setting $\theta_1>0$ and $\gamma_1=\gamma_2=\theta_2=0$ in the general fundamental equation~\eqref{eq:gen:fundamental_equation} (Lemma~\ref{lem:gen:fundamental_equation}) annihilates the $\partial P/\partial y$ term outright---its coefficient $xy\bigl[\gamma_2(y-x)+\theta_2(y-1)\bigr]$ vanishes identically---and trims the $\partial P/\partial x$ coefficient $xy\bigl[\gamma_1(x-y)+\theta_1(x-1)\bigr]$ down to its surviving abandonment part $\theta_1\,xy(x-1)$. What remains is
\begin{multline}
    \left[ (\lambda_1+\lambda_2+\mu)xy - \mu y - \lambda_1 x^2 y - \lambda_2 x y^2 \right] P(x,y)
    + \theta_1\, xy(x-1)\,\frac{\partial}{\partial x} P(x,y) \\
    \;=\; \mu (x-y)\,P_y(y) - \mu x(1-y)\,\pi(0,0),
    \label{eq:C2:remove}
\end{multline}
where, as in Model-$A$, the boundary function $P_y(y)=P(0,y)=\sum_{n_2\ge0} \pi(0,n_2)\,y^{n_2}$ is the partial PGF over the states with the server busy and the priority queue empty ($N_1=0$).

For each fixed $y$, \eqref{eq:C2:fundamental} is a first-order linear ODE in $x$ with two unknowns: $P(x,y)$ on $[0,1]^2$ and its boundary trace $P_y(y)$. We first pin $P_y(y)$ down (\S\S\ref{ssec:C2:anal}--\ref{ssec:C2:prob}) and then recover $P(x,y)$ (\S\ref{ssec:C2:recovery}).

The probabilistic route needs one identity. By the \emph{law of total expectation}, splitting $P_y(y)=\mathbb{E}[\mathbf{1}\{\text{busy},N_1=0\}\,y^{N_2}]$ over the event $\{\text{busy},N_1=0\}$ factors it into the mass of that event and the conditional PGF of $N_2$ given it:
\begin{equation}
    P_y(y) \;=\; \mathbb{P}(\text{busy},\, N_1=0)\,\cdot\,\mathbb{E}\!\left[y^{N_2} \,\big|\, \text{busy},\, N_1=0\right].
    \label{eq:C2:Py_prob_decomp}
\end{equation}

\subsection{Analytical derivation of \texorpdfstring{$P_y(y)$}{Py(y)} via boundedness at \texorpdfstring{$x=1$}{x=1}}\label{ssec:C2:anal}

The surviving abandonment factor $(x-1)$ places the operative singularity on the boundary $x=1$---the boundedness-pinning regime anticipated in Remark~\ref{rem:gen:reading}. Dividing~\eqref{eq:C2:fundamental} by the coefficient $\theta_1 xy(x-1)$ of the derivative puts it in the standard form~\eqref{eq:ode:standard} of \S\ref{sssec:if},
\begin{align*}
    \frac{\partial P}{\partial x} \;+\; \underbrace{\frac{(\lambda_1+\lambda_2+\mu)x - \mu - \lambda_1 x^2 - \lambda_2 xy}{\theta_1\,x(x-1)}}_{p(x;\,y)}\,P
    \;=\; \underbrace{\frac{\mu(x-y)\,P_y(y) - \mu x(1-y)\,\pi(0,0)}{\theta_1\,xy(x-1)}}_{q(x;\,y)}.
\end{align*}
The numerator of $p(x;y)$ is the kernel polynomial $A(x,y)$ of Model-$B$~\eqref{eq:B:Adef}, common to every model, with roots $x^*(y),x^+(y)$~\eqref{eq:A:xstar} satisfying $x^*(y)\,x^+(y)=\mu/\lambda_1$. Its partial-fraction decomposition is
\begin{equation*}
    \frac{A(x,y)}{x(x-1)} \;=\; -\lambda_1 + \frac{B}{x} + \frac{C}{x-1},
\end{equation*}
with $-\lambda_1$ the ratio of leading ($x^2$) coefficients, $B=A(0,y)/(0-1)=(-\mu)/(-1)=\mu$ by the cover-up rule at $x=0$, and $C=A(1,y)=\lambda_2(1-y)$ by the cover-up rule at $x=1$ (since $A(1,y)=(\lambda_1+\lambda_2+\mu)-\mu-\lambda_1-\lambda_2 y=\lambda_2(1-y)$). Thus
\begin{equation}
    p(x;y) \;=\; \frac{1}{\theta_1}\!\left[-\lambda_1 + \frac{\mu}{x} + \frac{\lambda_2(1-y)}{x-1}\right].
    \label{eq:C2:pfd}
\end{equation}
Integrating term by term, $\int p\,dx = \theta_1^{-1}\bigl[-\lambda_1 x + \mu\ln x + \lambda_2(1-y)\ln(1-x)\bigr]$, which exponentiates to the integrating factor
\begin{equation}
    \mu_I(x;y) \;=\; e^{-\lambda_1 x/\theta_1}\, x^{\alpha}\, (1-x)^{\beta(y)},
    \label{eq:C2:integrating_factor}
\end{equation}
where, in parallel with Model-$B_2$, we set
\begin{equation}
    \alpha \;=\; \frac{\mu}{\theta_1}, \qquad \beta(y) \;=\; \frac{\lambda_2(1-y)}{\theta_1}.
    \label{eq:C2:exponents}
\end{equation}
Both exponents are non-negative ($\alpha>0$ always; $\beta(y)\ge0$ for $y\in[0,1]$), so $\mu_I$ vanishes at $x=0$ and at $x=1$. By the integrating-factor solution~\eqref{eq:ode:general} of \S\ref{sssec:if}---with the boundary term $P(0,y)\,\mu_I(0;y)$ killed by $\mu_I(0;y)=0$---we obtain
\begin{equation}
    P(x,y) \;=\; \frac{1}{\mu_I(x;y)} \int_0^x q(t;y)\,\mu_I(t;y)\,dt.
    \label{eq:C2:Pxy_implicit}
\end{equation}

We now pin $P_y(y)$ by the same boundedness argument used for Model-$A$, watching $P(x,y)$ as $x\to 1$ for fixed $y\in[0,1)$: there $\beta(y)>0$, so $\mu_I(x;y)\to0$. (At $y=1$ the exponent $\beta(1)=0$ and $\mu_I(\,\cdot\,;1)$ no longer vanishes at $x=1$; the value $P_y(1)$ is then recovered by continuity, $y\to1^-$.) Since the left-hand side $P(x,y)$ of~\eqref{eq:C2:Pxy_implicit} stays bounded while $1/\mu_I(x;y)\to\infty$, the integral must vanish in the limit:
\begin{equation}
    \int_0^1 q(t;y)\,\mu_I(t;y)\,dt \;=\; 0.
    \label{eq:C2:boundedness}
\end{equation}
Substituting the explicit $q$ and $\mu_I$, the integrand is
\begin{equation*}
    q(t;y)\,\mu_I(t;y)
    \;=\; \frac{\mu\bigl[(t-y)\,P_y(y) - t(1-y)\,\pi(0,0)\bigr]}{\theta_1\,y}
    \cdot \frac{e^{-\lambda_1 t/\theta_1}\,t^{\alpha}(1-t)^{\beta(y)}}{t\,(t-1)},
\end{equation*}
where $t^{\alpha}/t=t^{\alpha-1}$ and the elementary identity $(1-t)^{\beta(y)}/(t-1)=-(1-t)^{\beta(y)-1}$ flips the sign and lowers the exponent. The constant prefactor $-\mu/(\theta_1 y)$, nonzero for $y\in(0,1]$, divides out of~\eqref{eq:C2:boundedness}, leaving
\begin{equation*}
    \int_0^1 e^{-\lambda_1 t/\theta_1}\,t^{\,\alpha-1}(1-t)^{\,\beta(y)-1}\bigl[(t-y)\,P_y(y) - t(1-y)\,\pi(0,0)\bigr]\,dt \;=\; 0.
\end{equation*}
Introduce the auxiliary integrals
\begin{align*}
    \mathcal{J}_0(y) &\;=\; \int_0^1 e^{-\lambda_1 t/\theta_1}\,t^{\,\alpha-1}(1-t)^{\,\beta(y)-1}\,dt, \\
    \mathcal{J}_1(y) &\;=\; \int_0^1 e^{-\lambda_1 t/\theta_1}\,t^{\,\alpha}(1-t)^{\,\beta(y)-1}\,dt.
\end{align*}
Distributing the bracket over $t^{\alpha-1}$, the term $t\cdot t^{\alpha-1}=t^{\alpha}$ feeds $\mathcal{J}_1$ and $t^{\alpha-1}$ feeds $\mathcal{J}_0$, so the integral becomes
\begin{align*}
    0 \;&=\; P_y(y)\!\int_0^1\! e^{-\lambda_1 t/\theta_1}\bigl(t^{\alpha}-y\,t^{\alpha-1}\bigr)(1-t)^{\beta(y)-1}\,dt
    \;-\; (1-y)\,\pi(0,0)\!\int_0^1\! e^{-\lambda_1 t/\theta_1}\,t^{\alpha}(1-t)^{\beta(y)-1}\,dt \\
    \;&=\; P_y(y)\bigl[\mathcal{J}_1(y) - y\,\mathcal{J}_0(y)\bigr] \;-\; (1-y)\,\pi(0,0)\,\mathcal{J}_1(y),
\end{align*}
which gives
\begin{equation}
    P_y(y) \;=\; \pi(0,0)\,(1-y)\,\frac{\mathcal{J}_1(y)}{\mathcal{J}_1(y) - y\,\mathcal{J}_0(y)}.
    \label{eq:C2:Py_anal_raw}
\end{equation}

It remains to evaluate the ratio $\mathcal{J}_1/\mathcal{J}_0$. Setting $a=\alpha=\mu/\theta_1$, $b=s/\theta_1$, $c=\lambda_1/\theta_1$, the Euler representation~\eqref{eq:kummer:euler} writes each $\mathcal{J}_i$ as a Beta multiple of a $\,{}_1F_1$ value, and the closed-form summation~\eqref{eq:erlangA:cf_kummer} of \S\ref{ssec:erlangA} identifies their ratio with the busy-period LST:
\begin{equation}
    \widetilde{B}_C(s) \;=\; \frac{\displaystyle\int_{0}^{1} t^{a}(1-t)^{b-1}e^{-ct}\,dt}{\displaystyle\int_{0}^{1} t^{a-1}(1-t)^{b-1}e^{-ct}\,dt}
    \;=\; \frac{B(a+1,b)\,{}_1F_1(a+1;\,a+1+b;\,-c)}{B(a,b)\,{}_1F_1(a;\,a+b;\,-c)}.
    \label{eq:C2:Bc_Kummer}
\end{equation}
Taking $s=\lambda_2(1-y)$---so that $b=\beta(y)$ matches the exponent in~\eqref{eq:C2:exponents}---gives $\widetilde{B}_C(\lambda_2(1-y))=\mathcal{J}_1(y)/\mathcal{J}_0(y)$. Dividing the numerator and denominator of~\eqref{eq:C2:Py_anal_raw} by $\mathcal{J}_0(y)$ yields the boundary generating function
\begin{equation}
    P_y(y) \;=\; \pi(0,0)\,\frac{(1-y)\,\widetilde{B}_C(\lambda_2(1-y))}{\widetilde{B}_C(\lambda_2(1-y)) - y}.
    \label{eq:C2:Py}
\end{equation}
This route uses no probabilistic structure: it reaches~\eqref{eq:C2:Py} through the Beta-integral identity~\eqref{eq:C2:Bc_Kummer} alone---the analytic shadow of the $M/G/1$/busy-period decomposition that the next subsection makes explicit.

\subsection{Probabilistic re-derivation of \texorpdfstring{$P_y(y)$}{Py(y)}}\label{ssec:C2:prob}

We again take the point of view of a class-$2$ customer waiting while the server is busy. Unlike when jockeying, class-$1$ customers here truly leave the system, so the class-$2$ Poisson arrival stream is undisturbed. the class-$2$ \emph{effective service time} is therefore one full busy period $B_C$ of the class-$1$ subsystem —now an $M/M/1{+}M$ queue with rates $\lambda_1,\mu,\theta_1$, rather than the pure $M/M/1$ of Model-$A$. So class-$2$ sees an $M/G/1$ queue with Poisson input $\lambda_2$ and service distributed as $B_C$. Its LST $\widetilde{B}_C(s)$ and mean $\mathbb{E}[B_C]$ come from the first-step analysis of \S\ref{ssec:erlangA} (Equations~\eqref{eq:erlangA:bp_lst} and \eqref{eq:erlangA:bp_mean}).

This $M/G/1$ is stable if and only if its utilisation $\lambda_2\,\mathbb{E}[B_C] < 1$ (Equation~\eqref{eq:C2:stability}). The condition is also strictly weaker than the no-abandonment requirement $\rho=\rho_1+\rho_2<1$.

Applying the Pollaczek-Khinchine formula~\eqref{eq:pk:pgf} with the identifications $\lambda\mapsto\lambda_2$, $\widetilde{B}\mapsto\widetilde{B}_C$, $\rho_B=\lambda_2\,\mathbb{E}[B_C]$ gives the PGF of the stationary number of class-$2$ customers in that $M/G/1$:
\begin{equation*}
    L_2(y) \;=\; \frac{(1 - \lambda_2 \mathbb{E}[B_C])\,(1-y)\,\widetilde{B}_C(\lambda_2(1-y))}{\widetilde{B}_C(\lambda_2(1-y)) - y}.
\end{equation*}
Since class-$2$ arrivals are Poisson-distributed, by \emph{PASTA} (\S\ref{ssec:priority}), the stationary number of class-$2$ customers is that of the $M/G/1$ at all times. Conditioning on $\{\text{busy},\,N_1=0\}$, we can identify that with that count with $N_2$:
\begin{equation}
    \mathbb{E}\!\left[y^{N_2} \,\big|\, \text{busy},\, N_1=0\right] \;=\; L_2(y).
    \label{eq:C2:L2_identification}
\end{equation}
Substituting~\eqref{eq:C2:L2_identification} into the decomposition~\eqref{eq:C2:Py_prob_decomp} gives
\begin{equation}
    P_y(y) \;=\; \mathbb{P}(\text{busy},\,N_1=0)\,\frac{(1-\lambda_2\mathbb{E}[B_C])\,(1-y)\,\widetilde{B}_C(\lambda_2(1-y))}{\widetilde{B}_C(\lambda_2(1-y)) - y}.
    \label{eq:C2:Py_unsimplified}
\end{equation}
The unknown mass $\mathbb{P}(\text{busy},N_1=0)$ is fixed by $P_y(0)=\pi(0,0)$: evaluating~\eqref{eq:C2:Py_unsimplified} at $y=0$, where $\widetilde{B}_C(\lambda_2)$ cancels,
\begin{align*}
    \pi(0,0)=P_y(0)
    &= \mathbb{P}(\text{busy},\,N_1=0)\,\frac{(1-\lambda_2\mathbb{E}[B_C])\cdot 1\cdot \widetilde{B}_C(\lambda_2)}{\widetilde{B}_C(\lambda_2)}
    &= \mathbb{P}(\text{busy},\,N_1=0)\,(1-\lambda_2\mathbb{E}[B_C]),
\end{align*}
so $\mathbb{P}(\text{busy},N_1=0)=\pi(0,0)/(1-\lambda_2\mathbb{E}[B_C])$. Substituting back, the factor $(1-\lambda_2\mathbb{E}[B_C])$ cancels and we recover
\begin{equation*}
    P_y(y) \;=\; \pi(0,0)\,\frac{(1-y)\,\widetilde{B}_C(\lambda_2(1-y))}{\widetilde{B}_C(\lambda_2(1-y)) - y},
\end{equation*}
which matches the result of Equation~\eqref{eq:C2:Py}, and so does match the expression as in the probabilistic argument for Model-$A$ (Equation~\eqref{eq:A:Py}); the only difference is the effective service time, now the busy period of an $M/M/1{+}M$ (Erlang-A) subsystem rather than of an ordinary $M/M/1$.

\subsection{Recovery of \texorpdfstring{$P(x,y)$}{P(x,y)} from \texorpdfstring{$P_y(y)$}{Py(y)}}\label{ssec:C2:recovery}

With $P_y(y)$ in hand, we return to the ODE~\eqref{eq:C2:fundamental} and substitute~\eqref{eq:C2:Py} into its right-hand side. Writing $\zeta(y)\equiv\widetilde{B}_C(\lambda_2(1-y))$ for brevity,
\begin{align*}
    \text{RHS}
    &\;=\; \mu(x-y)\,\pi(0,0)\,\frac{(1-y)\,\zeta(y)}{\zeta(y) - y} \;-\; \mu x(1-y)\,\pi(0,0) \\
    &\;=\; \mu\,\pi(0,0)\,(1-y)\!\left[\frac{(x-y)\,\zeta(y)}{\zeta(y) - y} - x\right] \\
    &\;=\; \mu\,\pi(0,0)\,(1-y)\,\frac{(x-y)\,\zeta(y) - x\bigl[\zeta(y) - y\bigr]}{\zeta(y) - y} \\
    &\;=\; \mu\,\pi(0,0)\,y(1-y)\,\frac{x - \zeta(y)}{\zeta(y) - y},
\end{align*}
which has the same form as the RHS of Model-$A$, now with the more involved service time $B_C$. With $p(x;y)$ and $\mu_I(x;y)$ as in~\eqref{eq:C2:pfd} and~\eqref{eq:C2:integrating_factor}, the standard-form source becomes
\begin{equation*}
    q(x;y) \;=\; \frac{\mu\,\pi(0,0)\,(1-y)\,\bigl[x-\zeta(y)\bigr]}{\theta_1\,x(x-1)\,\bigl[\zeta(y)-y\bigr]},
\end{equation*}
and~\eqref{eq:C2:Pxy_implicit} reads $P(x,y)=\mu_I(x;y)^{-1}\int_0^x q(t;y)\,\mu_I(t;y)\,dt$. To carry out the integral, use the partial fraction
\begin{equation*}
    \frac{t-\zeta(y)}{t(t-1)}
    \;=\; \frac{\zeta(y)(t-1) + t\bigl(1-\zeta(y)\bigr)}{t(t-1)}
    \;=\; \frac{\zeta(y)}{t} + \frac{1-\zeta(y)}{t-1},
\end{equation*}
and define the auxiliary integrals
\begin{align*}
    \mathcal{H}_1(x,y) &\;=\; \int_0^x e^{-\lambda_1 t/\theta_1}\,t^{\,\alpha-1}(1-t)^{\,\beta(y)}\,dt,  \\
    \mathcal{H}_2(x,y) &\;=\; \int_0^x e^{-\lambda_1 t/\theta_1}\,t^{\,\alpha}(1-t)^{\,\beta(y)-1}\,dt,
\end{align*}
with $\alpha,\beta(y)$ as in~\eqref{eq:C2:exponents}. These are the partial (upper limit $x$) counterparts of the full integrals of \S\ref{ssec:C2:anal}: $\mathcal{H}_2(1,y)=\mathcal{J}_1(y)$ exactly, whereas $\mathcal{H}_1(1,y)\neq\mathcal{J}_0(y)$ (its $(1-t)$ exponent is $\beta(y)$, not $\beta(y)-1$). The mismatch is harmless, as $\mathcal{H}_1$ and $\mathcal{H}_2$ only ever appear in the combination $\zeta\,\mathcal{H}_1-(1-\zeta)\,\mathcal{H}_2$.

Multiplying $q(t;y)$ by $\mu_I(t;y)=e^{-\lambda_1 t/\theta_1}t^{\alpha}(1-t)^{\beta(y)}$ and inserting the partial fraction, the integrand splits into the two pieces defining $\mathcal{H}_1,\mathcal{H}_2$:
\begin{align*}
    q(t;y)\,\mu_I(t;y)
    &\;=\; \frac{\mu\,\pi(0,0)\,(1-y)}{\theta_1\bigl(\zeta(y)-y\bigr)}
      \left[\frac{\zeta(y)}{t} + \frac{1-\zeta(y)}{t-1}\right]
      e^{-\lambda_1 t/\theta_1}\,t^{\alpha}(1-t)^{\beta(y)} \\
    &\;=\; \frac{\mu\,\pi(0,0)\,(1-y)}{\theta_1\bigl(\zeta(y)-y\bigr)}
      \Bigl[\zeta(y)\,e^{-\lambda_1 t/\theta_1}\,t^{\alpha-1}(1-t)^{\beta(y)}
      - \bigl(1-\zeta(y)\bigr)\,e^{-\lambda_1 t/\theta_1}\,t^{\alpha}(1-t)^{\beta(y)-1}\Bigr],
\end{align*}
the second term again using $(1-t)^{\beta(y)}/(t-1)=-(1-t)^{\beta(y)-1}$, which produces both the minus sign and the exponent shift on $\mathcal{H}_2$. Since $\pi(0,0)$ and $\zeta(y)$ leave the integral,
\begin{equation*}
    \int_0^x q(t;y)\,\mu_I(t;y)\,dt \;=\; \frac{\mu\,\pi(0,0)\,(1-y)}{\theta_1\bigl(\zeta(y)-y\bigr)}\Bigl[\zeta(y)\,\mathcal{H}_1(x,y) \;-\; \bigl(1-\zeta(y)\bigr)\,\mathcal{H}_2(x,y)\Bigr].
\end{equation*}
Dividing by $\mu_I(x;y)$ and restoring $\zeta(y)=\widetilde{B}_C(\lambda_2(1-y))$, we arrive at the closed form
\begin{align}
    \begin{split}
        P(x,y) \;=\; \frac{e^{\lambda_1 x/\theta_1}\,\mu\,\pi(0,0)\,(1-y)}{\theta_1\, x^{\alpha}(1-x)^{\beta(y)}\bigl(\widetilde{B}_C(\lambda_2(1-y))-y\bigr)}\\[4pt]
        \times \Bigl[\widetilde{B}_C(\lambda_2(1-y))\,\mathcal{H}_1(x,y) - \bigl(1-\widetilde{B}_C(\lambda_2(1-y))\bigr)\,\mathcal{H}_2(x,y)\Bigr],
    \end{split}\label{eq:C2:Pxy}
\end{align}
which is the statement of the theorem: the boundary function $P_y(y)$---obtained analytically in \S\ref{ssec:C2:anal} and probabilistically in \S\ref{ssec:C2:prob}---combined with the analytic integration of the ODE in $x$.

\begin{proof}[Proof of Corollary~\ref{cor:C2:pi00}]
    The fully idle state---server empty, no customer present, probability $\pi_0$---communicates with the rest of the chain through a single neighbour: the busy-and-empty state $(N_1,N_2)=(0,0)$ of probability $\pi(0,0)$, with exactly one customer in service and both queues empty (recall the standing convention $\pi_0\neq\pi(0,0)$). The system leaves the idle state only on an arrival of either class, at combined rate $\lambda_1+\lambda_2$ (Poisson superposition), and re-enters it only on the service completion that empties the busy-and-empty state, at rate $\mu$ (no abandonment can fire from $\pi(0,0)$, as both queues are empty). Global balance across the cut isolating the idle state (as in~\eqref{eq:gen:balance_idle}) therefore gives
    \begin{equation}
        (\lambda_1+\lambda_2)\,\pi_0 \;=\; \mu\,\pi(0,0),
        \label{eq:C2:balance_zero}
    \end{equation}
    so $\pi(0,0)=(\rho_1+\rho_2)\,\pi_0$.

    To fix $\pi_0$, we again use the $M/G/1$ structure of \S\ref{ssec:C2:prob}. By the $M/G/1$ utilisation identity, the long-run fraction of time this queue is busy is $\rho_B=\lambda_2\,\mathbb{E}[B_C]$, and ``idle'' means no class-$2$ customer is present, so
    \begin{equation*}
        \mathbb{P}(\text{no class-$2$ present}) \;=\; 1 - \lambda_2\,\mathbb{E}[B_C].
    \end{equation*}
    Conditional on no class-$2$ being present, the priority subsystem evolves as the $M/M/1{+}M$ queue ($\lambda_1,\mu,\theta_1$) defining $B_C$, alternating idle periods and busy periods $B_C$ in an \emph{alternating renewal process} whose regenerations are the instants it empties. An idle period lasts until the next class-$1$ arrival, which by memorylessness is $\mathrm{Exp}(\lambda_1)$ with mean $1/\lambda_1$, so the mean cycle is $1/\lambda_1+\mathbb{E}[B_C]$. By the \emph{renewal--reward theorem} the idle fraction is the ratio of mean idle period to mean cycle:
    \begin{equation*}
        \mathbb{P}(\text{class-$1$ absent} \mid \text{no class-$2$ present})
        \;=\; \frac{1/\lambda_1}{1/\lambda_1+\mathbb{E}[B_C]}
        \;=\; \frac{1}{1+\lambda_1\,\mathbb{E}[B_C]}.
    \end{equation*}
    The fully idle event is the intersection of $A=\{\text{class-$1$ absent}\}$ and $B=\{\text{no class-$2$ present}\}$, so by the chain rule $\mathbb{P}(A\cap B)=\mathbb{P}(B)\,\mathbb{P}(A\mid B)$ (not independence),
    \begin{equation*}
        \pi_0 \;=\; \bigl(1-\lambda_2\,\mathbb{E}[B_C]\bigr)\cdot\frac{1}{1+\lambda_1\,\mathbb{E}[B_C]} \;=\; \frac{1-\lambda_2\,\mathbb{E}[B_C]}{1+\lambda_1\,\mathbb{E}[B_C]},
    \end{equation*}
    which is~\eqref{eq:C2:pi0}; substituting into~\eqref{eq:C2:balance_zero} gives~\eqref{eq:C2:pi00}.

    In the limit $\theta_1\to 0^+$ the $M/M/1{+}M$ busy period collapses to the pure-$M/M/1$ one, $\mathbb{E}[B_C]\to(\mu-\lambda_1)^{-1}$, so
    \[
        \lambda_1\,\mathbb{E}[B_C]\to\frac{\rho_1}{1-\rho_1},
        \qquad
        \lambda_2\,\mathbb{E}[B_C]\to\frac{\rho_2}{1-\rho_1},
    \]
    whence $1+\lambda_1\,\mathbb{E}[B_C]\to1/(1-\rho_1)$ and $1-\lambda_2\,\mathbb{E}[B_C]\to(1-\rho)/(1-\rho_1)$. Therefore $\pi_0\to1-\rho$ by~\eqref{eq:C2:pi0}, and $\pi(0,0)=(\rho_1+\rho_2)\,\pi_0\to\rho(1-\rho)$, recovering the Model-$A$ baseline.
\end{proof}

\subsection{Mean queue lengths and mean waiting times}

\subsubsection*{Why the diagonal no longer simplifies}

Setting $x=y=z$ in~\eqref{eq:C2:fundamental}, the abandonment derivative term $\theta_1 xy(x-1)\,\partial_xP$ survives (the factor $(x-1)|_{x=z}\neq0$ for $z\in(0,1)$), and the diagonal equation reads
\begin{equation}
    \bigl[(\lambda_1{+}\lambda_2{+}\mu)z^2-\mu z-\lambda z^3\bigr]P(z,z)
    + \theta_1 z^2(z-1)\,\frac{\partial P}{\partial x}(z,z)
    \;=\; -\mu z(1-z)\,\pi(0,0).
    \label{eq:C2:diag}
\end{equation}
The contrast with Model-$A$ —or models without abandonment— is instructive: $P(z,z)$ is no longer purely algebraic with closed form $P(z,z)=\pi(0,0)/(1-\rho z)$~\eqref{eq:A:Pzz}, where the total number of customers $\mathbb{E}[N_1]+\mathbb{E}[N_2]$ followed easily by differentiation. Here the surviving derivative term remains as a second unknown so that path is no longer available since abandonment removes customers and lowers both means relative to models without abandonments.

\subsubsection*{Class-2 mean queue length via \texorpdfstring{$P(1,y)$}{P(1,y)}}

Setting $x=1$ in~\eqref{eq:C2:fundamental} we obtain
\begin{equation*}
    \lambda_2 y(1-y)\,P(1,y) \;=\; \mu(1-y)\bigl[P_y(y)-\pi(0,0)\bigr]. \\
\end{equation*}
Cancelling $(1-y)$ for $y\neq1$ and substituting the known $P_y(y)$~\eqref{eq:C2:Py},
\begin{equation*}
    P(1,y)
    \;=\; \frac{\mu\bigl[P_y(y)-\pi(0,0)\bigr]}{\lambda_2 y}
    \;=\; \frac{\mu\,\pi(0,0)}{\lambda_2}\cdot\frac{1-\widetilde{B}_C(\lambda_2(1-y))}{\widetilde{B}_C(\lambda_2(1-y))-y},
\end{equation*}
the marginal PGF of $N_2$. It is $0/0$ at $y=1$, so set $\varepsilon=1-y\to0$ and use the LST expansion $\widetilde{B}_C(\lambda_2\varepsilon)=1-\lambda_2\mathbb{E}[B_C]\,\varepsilon+\tfrac{1}{2}\lambda_2^2\mathbb{E}[B_C^2]\,\varepsilon^2+O(\varepsilon^3)$:
\begin{align*}
    1-\widetilde{B}_C(\lambda_2\varepsilon)
    &\;=\; \lambda_2\mathbb{E}[B_C]\,\varepsilon
           -\tfrac{1}{2}\lambda_2^2\mathbb{E}[B_C^2]\,\varepsilon^2+O(\varepsilon^3), \\
    \widetilde{B}_C(\lambda_2\varepsilon)-y
    &\;=\; \bigl(1-\lambda_2\mathbb{E}[B_C]\,\varepsilon+\tfrac{1}{2}\lambda_2^2\mathbb{E}[B_C^2]\,\varepsilon^2\bigr)-(1-\varepsilon)
    \;=\; (1-\lambda_2\mathbb{E}[B_C])\,\varepsilon
           +\tfrac{1}{2}\lambda_2^2\mathbb{E}[B_C^2]\,\varepsilon^2+O(\varepsilon^3).
\end{align*}
Writing $a=\lambda_2\mathbb{E}[B_C]$ and $c=\tfrac{1}{2}\lambda_2^2\mathbb{E}[B_C^2]$,
\begin{equation*}
    P(1,1-\varepsilon)
    \;=\; \frac{\mu\,\pi(0,0)}{\lambda_2}\cdot\frac{a\varepsilon - c\varepsilon^2+O(\varepsilon^3)}{(1-a)\varepsilon+c\varepsilon^2+O(\varepsilon^3)}
    \;=\; \frac{\mu\,\pi(0,0)}{\lambda_2}\cdot\frac{a-c\varepsilon}{(1-a)+c\varepsilon}+O(\varepsilon^2).
\end{equation*}
Differentiating $f(\varepsilon)=(a-c\varepsilon)/((1-a)+c\varepsilon)$ by the quotient rule,
\begin{equation*}
    f'(\varepsilon)
    \;=\; \frac{-c\,\bigl((1-a)+c\varepsilon\bigr)-(a-c\varepsilon)\,c}{\bigl((1-a)+c\varepsilon\bigr)^2}
    \;=\; \frac{-c}{(1-a+c\varepsilon)^2},
\end{equation*}
and since $\tfrac{d}{dy}P(1,y)\big|_{y=1}=-\tfrac{d}{d\varepsilon}P(1,1-\varepsilon)\big|_{\varepsilon=0}$,
\begin{equation}
    \mathbb{E}[N_2]
    \;=\; \frac{d}{dy}P(1,y)\bigg|_{y=1}
    \;=\; \frac{\mu\,\pi(0,0)}{\lambda_2}\cdot\frac{c}{(1-a)^2}
    \;=\; \frac{\mu\,\pi(0,0)\,\lambda_2\,\mathbb{E}[B_C^2]}{2\,(1-\lambda_2\mathbb{E}[B_C])^2}.
    \label{eq:C2:EN2}
\end{equation}
As $\theta_1\to0^+$, $\mathbb{E}[B_C]\to(\mu-\lambda_1)^{-1}$ and $\mathbb{E}[B_C^2]\to2\mu/(\mu-\lambda_1)^3$, recovering the Model-$A$ value $\mathbb{E}[N_2]=\rho\rho_2/((1-\rho_1)(1-\rho))$.

\subsubsection*{Class-1 mean queue length}

We extract $\mathbb{E}[N_1]=\tfrac{\partial P}{\partial x}\big|_{(1,1)}$ by evaluating the theorem formula at $y=1$ and then differentiating in $x$.

\medskip\noindent\textit{Step~1: marginal PGF $P(x,1)$.}\enspace
As $y\to1$, the exponent $\beta(y)\to0$, so $(1-x)^{\beta(y)}\to1$ for fixed $x\in(0,1)$, and $\mathcal{H}_1(x,y)\to\int_0^x e^{-\lambda_1 t/\theta_1}\,t^{\,\mu/\theta_1-1}\,dt$ while $\mathcal{H}_2(x,y)$ tends to a finite limit. In the theorem's bracket the two integrals carry different weights: $\widetilde{B}_C(\lambda_2(1-y))\to1$, while $1-\widetilde{B}_C(\lambda_2(1-y))=\lambda_2\mathbb{E}[B_C]\,(1-y)+O((1-y)^2)\to0$, so the $\mathcal{H}_2$ term drops and the bracket collapses to $\mathcal{H}_1(x,1)$. The prefactor's $y$-dependent part satisfies $(1-y)/(\widetilde{B}_C(\lambda_2(1-y))-y)\to1/(1-\lambda_2\mathbb{E}[B_C])$, so
\begin{equation*}
    P(x,1)
    \;=\; \frac{\mu\,\pi(0,0)}{\theta_1\,(1-\lambda_2\mathbb{E}[B_C])}\,e^{\lambda_1 x/\theta_1}\,x^{-\mu/\theta_1}
    \int_0^x e^{-\lambda_1 t/\theta_1}\,t^{\,\mu/\theta_1-1}\,dt.
\end{equation*}

\medskip\noindent\textit{Step~2: differentiate $P(x,1)$ at $x=1$.}\enspace
Write $P(x,1)=C\,f(x)\,\mathcal{K}(x)$ with $C=\mu\pi(0,0)/(\theta_1(1-\lambda_2\mathbb{E}[B_C]))$, $f(x)=e^{\lambda_1 x/\theta_1}x^{-\mu/\theta_1}$ and $\mathcal{K}(x)=\int_0^x e^{-\lambda_1 t/\theta_1}t^{\,\mu/\theta_1-1}\,dt$. By the product rule with Leibniz's rule ($\mathcal{K}'(x)=e^{-\lambda_1 x/\theta_1}x^{\mu/\theta_1-1}$) and the logarithmic derivative $f'(x)/f(x)=(\lambda_1 x-\mu)/(\theta_1 x)$,
\begin{equation*}
    \frac{d}{dx}P(x,1)
    \;=\; C\Bigl[f'(x)\,\mathcal{K}(x)+f(x)\,e^{-\lambda_1 x/\theta_1}x^{\mu/\theta_1-1}\Bigr]
    \;=\; C\left[e^{\lambda_1 x/\theta_1}x^{-\mu/\theta_1}\,\frac{\lambda_1 x-\mu}{\theta_1 x}\,\mathcal{K}(x) + x^{-1}\right].
\end{equation*}

\medskip\noindent\textit{Step~3: evaluate at $x=1$.}\enspace
\begin{equation*}
    \mathbb{E}[N_1]
    \;=\; \left.\frac{d}{dx}P(x,1)\right|_{x=1}
    \;=\; \frac{\mu\,\pi(0,0)}{\theta_1\,(1-\lambda_2\mathbb{E}[B_C])}
    \left[e^{\lambda_1/\theta_1}\,\frac{\lambda_1-\mu}{\theta_1}\,\mathcal{K}(1) + 1\right],
\end{equation*}
where $\mathcal{K}(1)=\int_0^1 e^{-\lambda_1 t/\theta_1}t^{\,\mu/\theta_1-1}\,dt=(\theta_1/\lambda_1)^{\mu/\theta_1}\,\gamma(\mu/\theta_1,\,\lambda_1/\theta_1)$, with $\gamma(\cdot,\cdot)$ the lower incomplete gamma function. Since $\lambda_1-\mu<0$ ($\rho_1<1$) the bracket is positive. No elementary simplification is known, so $\mathbb{E}[N_1]$ is evaluated numerically; as $\theta_1\to0^+$ it recovers the Model-$A$ value~\eqref{eq:A:EN1}.

\subsubsection*{Mean waiting times via Little's Law}

The relation $\mathbb{E}[N_i]=\lambda_i\,\mathbb{E}[W_i]$ holds for any stable system. Here $\lambda_i$ is the arrival rate of \emph{all} class-$i$ customers, including those that eventually abandon, so $\mathbb{E}[W_i]$ is the mean queueing time averaged over both served and abandoned customers:
\begin{equation*}
    \mathbb{E}[W_1] \;=\; \frac{\mathbb{E}[N_1]}{\lambda_1},
    \qquad
    \mathbb{E}[W_2] \;=\; \frac{\mathbb{E}[N_2]}{\lambda_2}.
\end{equation*}
Using the throughput $\mu(1-\pi_0)$ instead would give the conditional waiting time of served customers only---a different object. As $\theta_1\to0^+$ both means recover the Model-$A$ values~\eqref{eq:A:EN1}--\eqref{eq:A:EW2}.
